% !TEX root = userguide.tex

\def\headdate{21st June 2012}

\section{Projection}
\label{sec:projection}

In this section examples for the projection add-on (see
Sec.~\ref{sec:projection_addon}) are given, including:
\begin{enumerate}
   \item projection of a vector defined on one mesh to a vector defined on
         another mesh,
   \item projection of a vector to another vector, both defined on the same
         mesh, but having a different shape function.
\end{enumerate}

\subsection{Definition}

The projection of field $u_1(\vek x)$, defined on mesh M1 using shape functions
$\phi_1(\vek x)$: $u_1(\vek x)=\sum_k u_1^k\phi_1^k(\vek x)$ on a space 
defined on mesh M2 using shape function
$\phi_2(\vek x)$: $u_2(\vek x)=\sum_k u_2^k\phi_2^k(\vek x)$ is given by
\begin{equation}
  \sum_m (\phi^k_2,\phi^m_2) u_2^m = (\phi^k_2,u_1)
\end{equation}
The field $u_2$ is basically a least-squares ``fit'' of $u_1$. The add-on
projection deals with the determination of the inner product of the
right-hand side. This involves finding the values of $u_1$ defined on mesh M1
in the integration points defined on mesh M2.

\subsection{Examples}

The example files are \texttt{projection1.f90} to 
\texttt{projection5.f90} 
and can be obtained by typing at the
command line:
\medskip
\begin{quote}
   \texttt{getexample projection}
\end{quote}
Please read the \texttt{README} file and the sources for more information
on the examples.
